% ============================================================
%  sections/00_abstract.tex
%  Abstract — ~150 từ
% ============================================================

\begin{abstract}
Sự phổ biến của các mô hình ngôn ngữ lớn (LLM) đã tạo ra làn sóng
văn bản do máy tạo ra (MGT) tràn lan trên các nền tảng mạng xã hội
đa ngôn ngữ, đặt ra thách thức cấp thiết về phát hiện và kiểm duyệt
nội dung. Các phương pháp hiện tại chủ yếu được thiết kế cho tiếng Anh
và đòi hỏi phần cứng GPU cao cấp, gây khó khăn cho việc triển khai
thực tế. Bài báo này trình bày một nghiên cứu so sánh hệ thống bốn
phương pháp phát hiện MGT trên dữ liệu mạng xã hội đa ngôn ngữ gồm
bốn ngôn ngữ: tiếng Anh (EN), tiếng Việt (VI), tiếng Trung (ZH) và
tiếng Ả Rập (AR). Chúng tôi xây dựng một bộ dữ liệu cân bằng gồm
15.986 bài đăng từ các corpus MXH thực tế, sau đó lần lượt đánh giá:
(1) bộ phát hiện tiền huấn luyện zero-shot, (2) tinh chỉnh XLM-RoBERTa
đa ngôn ngữ, (3) phát hiện thống kê dựa trên Perplexity, và (4) học
chuyển giao đa ngôn ngữ (leave-one-language-out). Phương
pháp tinh chỉnh XLM-RoBERTa đạt độ chính xác trung bình \textbf{81,4\%}
trên cả bốn ngôn ngữ, được huấn luyện trong 7,5 phút trên GPU
T4 với chi phí ước tính \textbf{\$1,50}. Kết quả chứng minh rằng phát
hiện MGT đa ngôn ngữ chất lượng cao hoàn toàn khả thi trên phần cứng
phổ thông, đồng thời cho thấy học chuyển giao liên ngôn ngữ là chưa
đủ, đặc biệt với tiếng Việt (F1 = 0,134 khi thiếu dữ liệu huấn luyện).

\keywords{Phát hiện văn bản do máy tạo \and Mạng xã hội đa ngôn ngữ
\and XLM-RoBERTa \and Perplexity \and Tiếng Việt \and Tiếng Trung \and
Tiếng Ả Rập}
\end{abstract}
